
\begin{tcolorbox}[
    enhanced,
    colback=white,
    colframe=black,
    arc=5mm,
    boxrule=1pt,
    left=2mm,
    right=2mm,
    top=1mm,
    bottom=1mm
]

\textbf{Role.} You simulate an end user who is searching a knowledge base.

\textbf{Input.} You receive a numbered list of passages; each item is prefixed with a \texttt{[Document ID]} tag.

\textbf{Task.} Pick exactly one passage as the answer source and produce a single natural-language question that satisfies all of the following constraints:\\
(a) The question must be uniquely answerable by the passage you picked, and not answerable (or only partially answerable) by any other passage in the list.\\
(b) The question must sound like a real end-user utterance. Do NOT reveal the existence of a document, do NOT quote a title verbatim, and do NOT use phrases such as "according to the passage".\\
(c) Prepend the corresponding \texttt{[Document ID]} tag, followed by a single space, followed by the question.

\textbf{Output format.} \texttt{[Document ID] <your question>}

\textbf{Few-shot demonstration.}\\
\{SEED\}

\textbf{Now process the following real input.}\\

\textbf{Passages:}\\
\{DOCUMENTS\}\\

\end{tcolorbox}
